\section{Agents as state machines}
\label{sec:efsm}

\subsection{Why a state machine describes an agent}
\label{sec:efsm-why}

A state machine is a simplified description of an agent. It omits the prompts, the model and the
tool code so that we can examine execution without describing the whole program.

An agent runs in discrete steps, such as model calls, tool calls and checks. Between steps, it
retains information and may receive new inputs. The program determines which data changes and
continuations are permitted. We can therefore describe the agent as a state machine with i) a set
of states, ii) inputs received during a step, iii) outputs emitted by a step, iv) an initial
configuration, and v) a transition rule that maps the current state and an input to a next state
and an output. A workflow intended to terminate also specifies terminal configurations, from
which no further step is scheduled.

\begin{definition}[State]
\label{def:state}
State: ``condition of the system at a particular point in time''. It ``[e]ncodes everything about
the past that influences the system's reaction to current input''~\citep[slide~10]{seshia2015discrete}.
\end{definition}

Here, state contains the information the program retains between steps. It needs to include only
the past information that affects the next step. A controller's next action depends on its
instructions and earlier observations. The runtime's next decision may depend on how many model
calls have been used. When the program retains both, two visits to the same controller can have
different successors: the code is the same, but the data it reads differ.

The system boundary determines what counts as an input. Model replies and tool results arrive
during a step, and the program cannot compute them from its retained data. We can represent
workspace contents in two ways: include the files in the modelled environment, or treat a tool's
report about them as an input observation. Under either choice, \textbf{the record determines the
permitted transitions. It does not determine which result arrives}.

\begin{definition}[Determinism]
\label{def:determinism}
Determinism: ``In every state, for all input values, exactly one (possibly implicit) transition is
enabled''~\citep[slide~29]{seshia2015discrete}.
\end{definition}

Whether the agent is deterministic depends on the system boundary. If each model reply is an
input, routing is deterministic whenever every routing rule is a function of the merged record.
The record enables exactly one outgoing edge, so the same record and reply always select the same
edge. If the model is inside the system and its reply is an unobserved internal choice, one record
can lead to either a tool call or a submission. This description is nondeterministic. It uses
nondeterminism in the standard way to represent unknown aspects of a system or its environment.

Neither reading assigns a probability to a reply. A probabilistic machine attaches a probability to
each transition, whereas for a nondeterministic one ``no such probability is known. We just know that
any of the enabled transitions from a state can be taken''~\citep[slides~33 and~37]{seshia2015discrete}.
Our descriptions identify possible continuations without assigning their likelihoods.
Deterministic routing alone does not make a run reproducible. Replay requires the same inputs,
including model replies that we do not control.

\subsection{Extended state machines}
\label{sec:efsm-extended}

An agent has a small number of control locations and a large amount of data. In a plain state
machine, each distinct conversation and counter value would require a separate state. Such a
diagram would make the control structure difficult to see.

\begin{definition}[Extended state machine]
\label{def:efsm}
``Extended state machines augment the FSM model with variables that may be read or
written''~\citep[slide~4]{seshia2015extended}.
\end{definition}

\begin{definition}[Configuration]
\label{def:configuration}
A configuration of an extended state machine is a pair $(q, v)$: the current control location $q$,
shown as a node in the diagram, and a valuation $v$ that assigns a value to every declared variable.
\end{definition}

An agent's control locations name its computations: plan, call the controller, run the tools,
check a candidate, and end. Its variables hold the conversation, observations, plan, counters and
status. In our agent, \demoref{(controller, model\_calls=2, \ldots)} and
\demoref{(controller, model\_calls=5, \ldots)} are two configurations at one location. Both run the
same code with different data. Drawing a separate node for each conversation would obscure this
shared control structure.

Finite control does not imply a finite configuration space. An extended state machine for a
traffic light at a pedestrian crossing has four locations and one variable, \emph{count}, whose
range is $\{0, \ldots, 60\}$. It therefore has at most $4 \times 61 = 244$ configurations. A message
list without a length bound has no such limit. \textbf{The control locations are finite; the
configurations need not be.} Enumerating reachable configurations, as in safety analysis of a
finite machine, may then no longer terminate. The set of locations and the guards between
them remains finite.

\subsection{Transition notation}
\label{sec:efsm-notation}

An edge of an extended state machine can specify a guard, an output action and a set action. The
guard determines whether the transition is enabled; the output action specifies what it emits;
the set action updates the variables. We write \emph{guard / output action} above the edge and the
set action below it. For example, the traffic-light controller's red-to-green edge has
$\mathit{count} \geq 60 \,/\, \mathit{sigG}$ above $\mathit{count} := 0$.
The initial arrow carries the set action that initializes the variables. We declare variables,
inputs and outputs separately, following the notation's convention~\citep[slides~5--6]{seshia2015extended}: ``We make explicit declarations of variables, inputs,
and outputs to help distinguish the three.''

In prose, we use a one-line shorthand. A transition from configuration $(q, v)$ reads input $e$,
tests guard $g$, emits output \emph{out} and applies set action \emph{set}:
\[
  (q, v) \xrightarrow{\;e;\ g\ /\ \mathit{out};\ \mathit{set}\;} (q', v').
\]
The guard is a predicate over the variables and input, $g(v, e)$. The set action changes $v$ to
$v'$. \Cref{tab:efsm-elements} explains each element for an agent.

\begin{table}[htbp]
  \centering
  \begin{adjustbox}{width=\linewidth}\begin{tabularx}{\linewidth}{@{}L{0.24\linewidth}L{0.15\linewidth}X@{}}
    \toprule
    \textbf{Element} & \textbf{Symbol} & \textbf{In an agent} \\
    \midrule
    Control location & $q$ & The computation ready to run: a model call, a tool dispatch, a check.
    There are finitely many, each shown as one node. \\
    Variables & $v$ & What the program retains between steps: messages, observations, a plan,
    counters, check results, a status. \\
    Input & $e$ & A result that arrives during a step: a model reply, a tool result, a check
    result. \\
    Guard & $g(v, e)$ & A predicate over variables and input, such as \emph{the reply requests a tool} or
    \emph{the call budget is spent}. \\
    Output action & \emph{out} & A request to perform an operation, or a reported terminal
    outcome. \\
    Set action & \emph{set} & Store a reply, extend the observations, increment a counter, replace a
    plan. \\
    Initial configuration & $(q_0, v_0)$ & The first location, with the initial set action: opening
    messages present, counters at zero, no observations, status running. \\
    Terminal configuration & & A configuration from which no further computation is scheduled; the
    record carries the reason for ending. \\
    \bottomrule
  \end{tabularx}\end{adjustbox}
  \caption{The elements of an extended state machine, interpreted for an agent. Adapted from the general
  notation for extended state machines.}
  \label{tab:efsm-elements}
\end{table}

Consider a transition from the controller. The call counter $c$ was below the call budget $B$, so
the controller could issue a model call. Its reply $e$ contains a tool request. The reply is the
input, the guard tests for tool calls, and the output action requests their execution. The set
action appends the reply to the messages $H$ and increments $c$. The next location is \emph{tools}:
\[
  (\mathit{controller}, v) \xrightarrow{\;e;\ \mathit{calls}(e) \neq \emptyset\ /\ \mathit{request}(e);\
  H := H \cdot e,\ c := c + 1\;} (\mathit{tools}, v').
\]
When the tool results arrive, a transition from \emph{tools} appends them as observations and can
return control to the controller. This is the same location as before, with a different
configuration: $H$, $c$ and the observations have grown. Updating the variables allows the machine
to make progress while revisiting a location.

\subsection{Behaviours, traces and computation trees}
\label{sec:efsm-traces}

\begin{figure*}[tbp]
  \centering
  \begin{adjustbox}{width=\linewidth}% The planner as an extended state machine, labelled in the notation of sec:efsm-notation.
% Notes-owned copy (fig:efsm-planner); figures/lecture-3/plan-validate-efsm.tikz stays unchanged.
% Each label is guard / set action; output actions are omitted. Source: plan_and_validation/graph.py.
% The symbols are defined in a table set under the drawing in 07-efsm.tex (fig:efsm-planner).
\begin{tikzpicture}[
  ctl/.style={mnode, minimum width=21mm},
  run/.style={rnode, minimum width=21mm},
  el/.style={lab, font=\scriptsize, align=left},
  every node/.append style={text=ibmfg}
]
  \node[ctl] (plan) at (0,0) {plan};
  \node[ctl] (controller) at (6.4,0) {controller};
  \node[run] (tools) at (12.8,0) {tools};
  \node[run] (submit) at (6.4,-4.4) {submit};
  \node[ctl] (validate) at (12.8,-4.4) {validate};
  \node[tnode, minimum width=14mm] (end) at (0,-4.4) {END};

  % initial arrow with the initial set action
  \draw[flow] (-4.2,0) -- node[el, above, pos=.4]
    {$\sigma := \mathsf{running},\ c := 0,\ r := 0$\\
     $f := \emptyset,\ O := \langle\rangle,\ H := H_0$} (plan.west);

  % plan
  \draw[flow] (plan) -- node[el, above]
    {$c < B$ /\\ $P := e,\ f := \emptyset,\ c := c + 1$\\ $r := |O|$ if $f \neq \emptyset$} (controller);
  \draw[flow] (plan) -- node[el, left]
    {$c \geq B$ /\\ $\sigma := \mathsf{budget\_exhausted}$} (end);

  % controller
  \draw[flow] (controller) -- node[el, below]
    {$c < B \wedge \neg s \wedge \mathit{calls}(e) \neq \emptyset$ /\\ $H := H \cdot e,\ c := c + 1$} (tools);
  \draw[flow] (controller) -- node[el, right, pos=.55]
    {$c < B \wedge \neg s \wedge \mathit{calls}(e) = \emptyset$ /\\ $H := H \cdot e,\ c := c + 1$} (submit);
  \draw[flow] (controller.south west) -- node[el, below right, pos=.85]
    {$c \geq B$ / $\sigma := \mathsf{budget\_exhausted}$\\ $c < B \wedge s$ / $\sigma := \mathsf{no\_progress}$} (end.north east);

  % tools
  \draw[back] ([xshift=-3mm]tools.north) -- ++(0,1.3) -|
    node[el, above, pos=.25] {$\neg(r = 0 \wedge s)$ / $O := O \cdot e$} ([xshift=3mm]controller.north);
  \draw[back] ([xshift=3mm]tools.north) -- ++(0,2.5) -|
    node[el, above, pos=.25] {$r = 0 \wedge s$ / $O := O \cdot e,\ f := $ reason} (plan.north);

  % submit
  \draw[flow] (submit) -- node[el, above] {$\mathit{pass}(e)$ /\\ $K := e$} (validate);
  \draw[back] ([xshift=-5mm]submit.north) -- node[el, left, pos=.3]
    {$\neg\mathit{pass}(e)$ /\\ $K := e,\ H := H \cdot$ feedback} ([xshift=-5mm]controller.south);

  % validate
  \draw[flow] (validate.south) -- ++(0,-.9) -| node[el, below, pos=.25]
    {$c < B$ / $\sigma := \mathit{verdict}(e),\ c := c + 1$\\ $c \geq B$ / $\sigma := \mathsf{budget\_exhausted}$} (end.south);

\end{tikzpicture}
\end{adjustbox}

  \medskip
  \begin{adjustbox}{width=\linewidth}\begin{tabularx}{\linewidth}{@{}L{0.09\linewidth}XL{0.09\linewidth}X@{}}
    \toprule
    \textbf{Symbol} & \textbf{Meaning} & \textbf{Symbol} & \textbf{Meaning} \\
    \midrule
    $e$ & The step's input: a reply, tool results, or check results. &
      $O$ & Observations. \\
    $\sigma$ & Status. & $P$ & Plan. \\
    $c$ & Model calls completed. & $K$ & Check results. \\
    $B$ & Call budget (40). & $f$ & Failed step (\texttt{failed\_step}). \\
    $H$ & Messages. & $r$ & Number of observations at the last revision (\texttt{replanned\_at}). \\
    $H_0$ & The opening messages. & $s$ & The last three observations since index $r$, including
      any just appended, are identical, and none is a successful edit. \\
    $\mathit{calls}(e)$ & The tool calls in reply $e$. & $\mathit{pass}(e)$ & All three checks pass. \\
    $\mathit{verdict}(e)$ & \textsf{accepted} or \textsf{rejected}. & & \\
    \bottomrule
  \end{tabularx}\end{adjustbox}
  \caption{The planning agent as an extended state machine, drawn from its \demoref{graph.py}. Each edge is
  labelled \emph{guard / set action} in the notation of \cref{sec:efsm-notation}; output actions are
  omitted. Labels are stacked where two transitions share an arrow. Blue locations call a model;
  grey locations run tools or checks. The table under the drawing defines every symbol on the edges.}
  \label{fig:efsm-planner}
\end{figure*}

A transition specifies one permitted step. A run is a sequence of these steps. Three definitions
describe what a run did or could have done~\citep[slide~8]{seshia2015extended}.

\begin{definition}[Behaviour]
\label{def:behaviour}
``FSM behavior is a sequence of (non-stuttering) steps''.
\end{definition}

\begin{definition}[Trace]
\label{def:trace}
``A trace is the record of inputs, states, and outputs in a behavior''.
\end{definition}

\begin{definition}[Computation tree]
\label{def:computation-tree}
``A computation tree is a graphical representation of all possible traces''.
\end{definition}

A stuttering step occurs when inputs are absent, changes no state and emits no output. A behaviour
omits these steps. Determinism determines whether a system has one behaviour or many. For a fixed
input sequence, ``A deterministic system exhibits a single
behavior'', while ``A non-deterministic system exhibits a set of
behaviors''~\citep[slides~29 and~36]{seshia2015discrete}. The computation tree represents that set,
branching wherever more than one transition is enabled.

A trajectory is one path through the computation tree, from the root to a terminal configuration
or to the point where the run stopped. The tree unfolds the machine: each visit to a location
appears as a new node. Four visits to the controller therefore produce four tree nodes, each with
its own configuration and possible continuations. The machine diagram shows the permitted
transitions. The tree shows the possible sequences of transitions from the initial configuration.

An agent's event log is a practical trace that can survive the process. During execution,
in-memory state holds the run's current information. When the process exits, unsaved state is
lost, while saved logs and workspace artifacts can remain. \textbf{State records what the run currently knows; the trace
records what the run did.} Our agent writes six kinds of event:
\demoref{tool\_request}, \demoref{tool\_result}, \demoref{model\_call}, \demoref{routing},
\demoref{state\_change} and \demoref{terminal}.

An event log meets \cref{def:trace} only partly. It contains the inputs and outputs selected for
logging and usually little state, because copying every variable at every step is expensive.
Reconstructing a trajectory therefore combines logged facts with inferences from transition
rules. A claim about a run needs to identify which of these supports it.

\subsection{The planning agent as an extended state machine}
\label{sec:efsm-demo}

Our running example is the \emph{planning agent}, the complete workflow. The \emph{planner} is its
plan node and model role. The agent receives one scanner alert and drafts a plan. A
reactive loop of controller and tool calls carries out the work. The agent submits a candidate
patch to three mechanical checks, then a validator decides whether to accept it. In
\cref{fig:efsm-planner}, one transition represents completing a computation, applying its set
action and selecting the next location. It omits the internal steps of model and tool calls. The
figures and runs describe the planning agent when both runs were recorded. The current code adds two
runtime nodes, \demoref{scan} and \demoref{normalize}, before \demoref{plan}.

The initial arrow enters \demoref{plan} with $\sigma := \mathsf{running}$, no model calls, no
observations, no revision ($r = 0$, $f = \emptyset$), and the opening system and task messages
$H_0$. \Cref{fig:efsm-planner} shows the following routes.

\bi

\item \textbf{Planning.} \demoref{plan} makes one model call, stores the proposed plan in $P$, and
moves to \demoref{controller}.

\item \textbf{Acting.} A controller reply with tool calls, $\mathit{calls}(e) \neq \emptyset$, routes
to \demoref{tools}, which appends the results to $O$. Unless the stall guard holds, control returns
to \demoref{controller}.

\item \textbf{Revising.} The stall guard is $r = 0 \wedge s$: the run has not revised its plan, and the
last three observations are identical with no successful edit among them. When it holds,
\demoref{tools} sets the handoff variable $f$ (\demoref{failed\_step}) and control returns to
\demoref{plan}, which replaces $P$, clears $f$, and sets $r := |O|$ (\demoref{replanned\_at}). From
then on, $s$ reads only the observations made since the revision.

\item \textbf{Submitting.} A reply without tool calls, $\mathit{calls}(e) = \emptyset$, routes to
\demoref{submit}, which runs the three checks: the patch applies, it touches the flagged line, and a
rescan is clean. On $\neg\mathit{pass}(e)$ it appends feedback to $H$ and returns to
\demoref{controller}; on $\mathit{pass}(e)$ it moves to \demoref{validate}.

\item \textbf{Judging.} \demoref{validate} makes one model call and always moves to \demoref{END},
with $\sigma := \mathit{verdict}(e)$, either \demoref{accepted} or \demoref{rejected}. A reply
without a decision is recorded as a rejection made by the runtime.

\item \textbf{Exhaustion.} \demoref{plan}, \demoref{controller} and \demoref{validate} each test
$c < B$ before calling a model, with $B = 40$; when the test fails, the edge to \demoref{END} sets
$\sigma := \mathsf{budget\_exhausted}$. At the controller, $s$ can hold only after a revision, and
it ends the run with $\sigma := \mathsf{no\_progress}$. The routing function after \demoref{submit}
also checks the status, but no step on the way to \demoref{submit} can change it.

\ei

The planner stores its list of subgoals in a variable. Adding a subgoal does not add a control
location, and routing does not test which subgoal is current. \textbf{A plan is data in the machine;
its subgoals do not define control locations.} The controller receives the plan as a message. The
plan affects the machine's behaviour only through the controller's replies.

Our worked run, \demoref{bc284c6b}, repairs a Bandit B608 alert (SQL injection) at line 46 of
\demoref{testcode/BenchmarkTest00283.py}. Its log records this trajectory, with the tool each
\demoref{tools} step ran:

\noindent\begin{minipage}{\linewidth}
\begin{lstlisting}[basicstyle=\ttfamily\scriptsize]
plan
  -> controller -> tools(read_file)
  -> controller -> tools(search)
  -> controller -> tools(read_file)
  -> controller -> tools(edit)
  -> controller -> submit -> validate -> END(accepted)
\end{lstlisting}
\end{minipage}

The trajectory contains 25 events. Events 1--2 record the model call in \demoref{plan} and the plan
it stored. Events 3--18 record four tool rounds. Each contains a controller call, routing decision,
tool request and tool result. The requests are \demoref{read\_file} on the flagged file,
\demoref{search} in \demoref{helpers/db\_sqlite.py}, \demoref{read\_file} on that helper, and
\demoref{edit} on the flagged file. Events 19--22 are submission: a controller reply with no tool
calls, the route to \demoref{submit}, three passed checks, and the route to \demoref{validate}.
Events 23--25 record validation and the terminal event, \demoref{accepted}. The run made seven
model calls: one by \demoref{plan}, five by \demoref{controller} and one by \demoref{validate}. It
logged 12,693 input and 3,755 output tokens.

\begin{figure}[tbp]
  \centering
  \begin{adjustbox}{width=\linewidth}% A fragment of the planner's computation tree around run bc284c6b (fig:efsm-trace-tree).
% Notes-owned; figures/lecture-3/plan-validate-trace-tree.tikz stays unchanged.
% Branches are drawn where a reply or a check result selects the edge. Blue is the recorded path.
\begin{tikzpicture}[
  obs/.style={gnode, draw=ibmblue, fill=ibmblue!7, line width=.8pt, minimum width=19mm},
  alt/.style={gnode, draw=black!40, text=black!60, fill=black!2, minimum width=19mm},
  hit/.style={-{Stealth[length=2mm]}, line width=1.1pt, draw=ibmblue},
  may/.style={-{Stealth[length=2mm]}, densely dotted, line width=.8pt, draw=black!50},
  el/.style={lab, font=\scriptsize},
  cnt/.style={font=\scriptsize, text=ibmblue, anchor=east}
]
  \node[obs] (plan) at (0,0) {plan};
  \node[obs] (c1) at (0,-1.4) {controller};
  \node[obs] (t1) at (0,-2.8) {tools\\\texttt{read\_file}};
  \node[alt] (s1) at (3.6,-2.8) {submit};
  \node[obs] (c2) at (0,-4.2) {controller};
  \node[obs] (t2) at (0,-5.6) {tools\\\texttt{search}};
  \node[alt] (s2) at (3.6,-5.6) {submit};
  \node[obs, text width=30mm, align=center] (dots) at (0,-7.0) {$\cdots$\\two rounds:\\\texttt{read\_file}, \texttt{edit}};
  \node[obs] (c5) at (0,-8.6) {controller};
  \node[alt] (t5) at (-3.6,-10.0) {tools};
  \node[alt] (more) at (-3.6,-11.4) {$\cdots$};
  \node[alt, text width=24mm, align=center] (exh) at (-3.6,-12.8) {END\\\textsf{budget\_exhausted}};
  \node[obs] (sub) at (0,-10.0) {submit};
  \node[alt] (back) at (3.6,-10.0) {controller};
  \node[obs] (val) at (0,-11.4) {validate};
  \node[obs] (acc) at (0,-12.8) {END\\\textsf{accepted}};
  \node[alt] (rej) at (3.6,-12.8) {END\\\textsf{rejected}};

  \draw[hit] (plan) -- (c1);
  \draw[hit] (c1) -- node[el, left] {$\mathit{calls}(e) \neq \emptyset$} (t1);
  \draw[may] (c1) -- node[el, above right] {$\mathit{calls}(e) = \emptyset$} (s1);
  \draw[hit] (t1) -- (c2);
  \draw[hit] (c2) -- node[el, left] {$\mathit{calls}(e) \neq \emptyset$} (t2);
  \draw[may] (c2) -- node[el, above right] {$\mathit{calls}(e) = \emptyset$} (s2);
  \draw[hit] (t2) -- (dots);
  \draw[hit] (dots) -- (c5);
  \draw[may] (c5) -- node[el, above left] {$\mathit{calls}(e) \neq \emptyset$} (t5);
  \draw[hit] (c5) -- node[el, left] {$\mathit{calls}(e) = \emptyset$} (sub);
  \draw[may] (t5) -- (more);
  \draw[may] (more) -- node[el, left] {$c \geq B$} (exh);
  \draw[may] (sub) -- node[el, above] {$\neg\mathit{pass}(e)$} (back);
  \draw[hit] (sub) -- node[el, left] {$\mathit{pass}(e)$} (val);
  \draw[hit] (val) -- (acc);
  \draw[may] (val) -- (rej);

  \node[cnt] at (-1.2,-1.4) {$c = 1$};
  \node[cnt] at (-1.2,-4.2) {$c = 2$};
  \node[cnt] at (-1.2,-8.6) {$c = 5$};
  \node[cnt] at (-1.2,-11.4) {$c = 6$};
\end{tikzpicture}
\end{adjustbox}
  \caption{A fragment of the planning agent's computation tree around run \demoref{bc284c6b}. Blue nodes
  and edges show the recorded trajectory; the $\cdots$ node omits two controller and tool rounds.
  Dotted grey branches are continuations the guards permit under other replies or check results;
  the grey $\cdots$ omits further rounds. $c$ is the number of model calls completed on reaching a
  node, derived from the log and the set actions. No probabilities are implied.}
  \label{fig:efsm-trace-tree}
\end{figure}

\Cref{fig:efsm-trace-tree} unfolds the beginning and end of this trajectory into a tree. Each visit
to \demoref{controller} creates a new node whose reply selects \demoref{tools} or \demoref{submit}.
At \demoref{submit}, the check result selects the branch; at \demoref{validate}, the verdict does.
An exhaustion branch appears only where its guard permits it. On the recorded path, $c$ never
exceeds 7, so reaching \demoref{END} with \demoref{budget\_exhausted} would require further rounds.
The figure shows branches selected by replies or check results. Tool results can also change the
route through the stall guard; those branches are omitted.

The log contains selected facts. A successful \demoref{tool\_result} records success and result
length, so the observations themselves are absent. The routes from \demoref{tools} to
\demoref{controller} and from \demoref{plan} to \demoref{controller} produce no routing event. We
infer these routes because the next event is a controller call and the transition rules allow no
other successor. The counters are also absent from the log. Reconstructing this trajectory
therefore requires both the log and the transition rules. Recovering every variable value at
every step would require more detailed records.

\subsection{Reading}

\citet{seshia2015discrete}, slides 10--37, on state, transitions and guards,
determinism and receptiveness, and nondeterministic machines and their computation trees; and
\citet{seshia2015extended}, slides 4--8, on extended state machines, their notation, the
traffic-light example, and behaviours and traces.

\paragraph{Looking ahead.} The machine has fixed control locations and variables whose values
change at each step. The next section describes the record that holds those variables, the steps
that may update each field, and the rules for merging updates with stored values.
